%%%%%%%%%%%%%%%%% Add paper-specific macros here
\newcommand{\sysname}{LiVo\xspace}

%%% Latin abbrevs
\newcommand{\etc}{\emph{etc.}\xspace}
\newcommand{\ie}{\emph{i.e.,}\xspace}
\newcommand{\eg}{\emph{e.g.,}\xspace}
\newcommand{\etal}{\emph{et al.}\xspace}

%%% Sections, figures etc. Please use these for uniformity
\newcommand{\secref}[1]{\S\ref{#1}}
\newcommand{\figref}[1]{Fig.~\ref{#1}}
\newcommand{\tabref}[1]{Tbl.~\ref{#1}}
\newcommand{\algoref}[1]{Algorithm~\ref{#1}}
\newcommand{\eqnref}[1]{Equation~\ref{#1}}
\newcommand{\lineref}[1]{Line~\ref{#1}}
\newcommand{\algolineref}[2]{Alg.~\ref{#1}, Line~\ref{#2}}
\newcommand{\algolinesref}[3]{Alg.~\ref{#1}, Lines~\ref{#2}--\ref{#3}}
\newcommand{\custombullet}{\paraspace\noindent$\blacktriangleright$}

\newcommand\checkmarkk[1][]{%
  \tikz[scale=0.4,#1]{\fill(0,.35) -- (.25,0) -- (1,.7) -- (.25,.15) -- cycle;}%
}
\newcommand\crossmark[1][]{%
  \tikz[scale=0.4,#1]{
    \fill(0,0)--(0.1,0) .. controls (0.5,0.4) .. (1,0.7)--(0.9,0.7) ..  controls (0.5,0.5) ..(0,0.1) --cycle;
    \fill(1,0.1)--(0.9,0.1) .. controls (0.5,0.3) .. (0,0.7)--(0.1,0.7) .. controls (0.5,0.4) ..(1,0.2) --cycle;
  }%
}

\newcommand\ckmark[1]{\ding{52}}

%%% Commenting macros: change as necessary
\newcommand{\ramesh}[1]{\todo[author=Ramesh,color=blue!10,inline]{#1}}
\newcommand{\raj}[1]{\todo[author=Raj,color=cyan!25,inline]{#1}}
\newcommand{\christina}[1]{\todo[author=Christina,color=pink,inline]{#1}}
\newcommand{\rn}[1]{\todo[author=Ravi,color=red!20,inline]{#1}}
\newcommand{\harsha}[1]{\todo[author=Harsha,color=blue!10,inline]{#1}}
\newcommand{\antonio}[1]{\todo[author=Antonio,color=green!10,inline]{#1}}
\newcommand{\sanjay}[1]{\todo[author=Sanjay,color=green!20,inline]{#1}}
\newcommand{\tao}[1]{\todo[author=Tao,color=purple!10,inline]{#1}}
\newcommand{\agr}[1]{\todo[author=Anthony,color=orange!20,inline]{#1}}

\newcommand{\note}[1]{\textcolor{red}{#1}}

%%% Camera-ready macros
\newcommand{\cradd}[1]{\textcolor{blue}{#1}}
\newcommand{\crdel}[1]{\sout{\textcolor{red}{#1}}}

%%% For argmin
\DeclareMathOperator*{\argmin}{arg\,min}

%%% For tight lists
\newcommand{\squishlist}{
   \begin{list}{$\bullet$}
    { \setlength{\itemsep}{0pt}      \setlength{\parsep}{1pt}
      \setlength{\topsep}{1pt}       \setlength{\partopsep}{0pt}
      \setlength{\leftmargin}{1.0em} \setlength{\labelwidth}{1em}
      \setlength{\labelsep}{0.5em} } }

\newcommand{\squishlisttwo}{
   \begin{list}{$\bullet$}
    { \setlength{\itemsep}{0pt}    \setlength{\parsep}{0pt}
      \setlength{\topsep}{0pt}     \setlength{\partopsep}{0pt}
      \setlength{\leftmargin}{2em} \setlength{\labelwidth}{1.5em}
      \setlength{\labelsep}{0.5em} } }

\newcommand{\squishend}{
    \end{list}  }

%%% For colors
\definecolor{Gainsboro}{rgb}{0.86, 0.86, 0.86}
